\chapter{Áp Lực Từ Cha Mẹ}
\markboth{Áp Lực Từ Cha Mẹ}{Pressure from Parents}

\section{Áp Lực Mang Tên Yêu Thương}
\begin{truyen}{Áp Lực Mang Tên Yêu Thương}{Pressure That Arrives in the Name of Love}
\chuhoa{N}{am nghe câu này suốt ba năm cấp ba: ``Bố mẹ làm tất cả cũng vì con.''}
\textit{(Nam heard this sentence throughout high school: ``We do everything for you.'')}

Câu đó đúng một nửa.
\textit{(That sentence was half true.)}

Nửa còn lại là mọi kỳ vọng, mọi lịch học kín đặc, mọi so sánh ngầm cũng được gói trong chữ yêu.
\textit{(The other half was that every expectation, every overloaded schedule, every quiet comparison was also wrapped inside the language of love.)}

Nam không dám than mệt vì sợ bị gọi là vô ơn.
\textit{(Nam did not dare complain of exhaustion because he feared being called ungrateful.)}

Em không dám nói không thích ngành bố chọn vì sợ làm mẹ buồn.
\textit{(He did not dare say he disliked the major his father chose because he feared hurting his mother.)}

Một tối, em bật khóc và nói: "Con biết bố mẹ thương con.
\textit{(One night he cried and said: "I know you love me.)}

Nhưng nếu con chỉ được quyền làm người bố mẹ tự hào, còn con người thật của con thì sao?"
\textit{(But if I only get to be the child you are proud of, what happens to the real me?")}

Căn phòng im lặng rất lâu.
\textit{(The room stayed silent for a long time.)}

Đó là im lặng khó chịu, nhưng cần thiết.
\textit{(It was an uncomfortable silence, but a necessary one.)}

Có những gia đình chỉ bắt đầu chữa lành từ lúc một đứa trẻ dám nói thật.
\textit{(Some families only begin healing when a child dares to tell the truth.)}

\end{truyen}

\begin{baihoc}
Tình thương không nên biến thành chiếc lồng mạ vàng. Cha mẹ cần định hướng, nhưng cũng phải chừa chỗ cho tiếng nói thật của con.
\textit{(Love should not become a golden cage. Parents should guide, but must leave room for the child’s real voice.)}
\end{baihoc}

\begin{ghinhoanh}
\textbf{Short English to remember:}

Love should not become a golden cage.

Parents should guide, but must leave room for the child’s real voice.

\end{ghinhoanh}

\begin{tuhocnhanh}
\textbf{Cau hoi tu hoc / Self-study questions}

1. Van de chinh la gi? \textit{What is the main problem?}

2. Cach xu ly tot hon la gi? \textit{What is a better action?}

3. Mot cau tieng Anh em muon nho la cau nao? \textit{Which English sentence do you want to remember?}
\end{tuhocnhanh}
\ngancach

\section{So Với Con Nhà Người Ta Là Cách Giết Tự Tin Nhanh Nhất}
\begin{truyen}{So Với Con Nhà Người Ta Là Cách Giết Tự Tin Nhanh Nhất}{Comparing You to Other People’s Children Destroys Confidence Fastest}
\chuhoa{M}{ỗi lần Huy được tám điểm, bố hỏi: ``Sao không được mười như bạn An?''}
\textit{(Every time Huy scored an eight, his father asked: ``Why not ten like An?'')}

Mỗi lần Huy đá bóng giỏi, mẹ bảo: "Giỏi cái đó để làm gì?
\textit{(Whenever Huy played football well, his mother said: "What is the use of that?)}

Con nhìn con người ta học đi."
\textit{(Look at how other people’s children study.")}

Lâu dần, Huy không còn nghe lời góp ý.
\textit{(Over time, Huy no longer heard advice.)}

Em chỉ nghe thấy một thông điệp duy nhất: mình không bao giờ đủ.
\textit{(He heard only one message: I am never enough.)}

Người lớn thường tưởng so sánh sẽ tạo động lực.
\textit{(Adults often think comparison creates motivation.)}

Thực tế, so sánh kéo học sinh vào một cuộc đua mà đích đến luôn nằm ở người khác.
\textit{(In reality, comparison drags students into a race whose finish line is always inside someone else.)}

Khi một đứa trẻ mất cảm giác mình có giá trị riêng, nó hoặc lì ra, hoặc phát điên vì chứng minh, hoặc âm thầm ghét chính gia đình mình.
\textit{(When a child loses the sense of having unique worth, that child either goes numb, goes obsessive, or quietly grows to resent the family.)}

\end{truyen}

\begin{baihoc}
So sánh có thể tạo áp lực ngắn hạn, nhưng làm hỏng gốc rễ lâu dài. Hãy đo sự tiến bộ của một học sinh bằng chính đường đi của em ấy, không phải bằng con nhà người ta.
\textit{(Comparison may create short-term pressure, but it damages the roots long-term. Measure a student’s growth against that student’s own path, not somebody else’s child.)}
\end{baihoc}

\begin{ghinhoanh}
\textbf{Short English to remember:}

Comparison may create short-term pressure, but it damages the roots long-term.

Measure a student’s growth against that student’s own path, not somebody else’s child.

\end{ghinhoanh}

\begin{tuhocnhanh}
\textbf{Cau hoi tu hoc / Self-study questions}

1. Van de chinh la gi? \textit{What is the main problem?}

2. Cach xu ly tot hon la gi? \textit{What is a better action?}

3. Mot cau tieng Anh em muon nho la cau nao? \textit{Which English sentence do you want to remember?}
\end{tuhocnhanh}
\ngancach

\section{Đứa Con Ngoan Không Phải Lúc Nào Cũng Ổn}
\begin{truyen}{Đứa Con Ngoan Không Phải Lúc Nào Cũng Ổn}{The Good Child Is Not Always Fine}
\chuhoa{C}{ó những đứa trẻ không cãi, không phá, không gây rắc rối.}
\textit{(There are children who do not argue, do not rebel, and do not cause trouble.)}

Người lớn rất thích kiểu con như vậy.
\textit{(Adults love that kind of child.)}

Nhưng đôi khi sự ngoan đó không phải trưởng thành.
\textit{(But sometimes that obedience is not maturity.)}

Nó là kết quả của việc sợ làm ai đó thất vọng.
\textit{(It is the result of being terrified of disappointing someone.)}

Thu luôn nói: ``Dạ, sao cũng được.'' Em hiếm khi xin thứ mình muốn, hiếm khi phản đối điều mình không thích.
\textit{(Thu always said: ``Yes, whatever you think.'' She rarely asked for what she wanted, and rarely objected to what she disliked.)}

Đến khi vào đại học, không ai quản lịch nữa, em sụp hoàn toàn.
\textit{(When she entered university and no one managed her schedule anymore, she collapsed.)}

Em không biết mình thích gì, giỏi gì, chọn gì, vì cả tuổi học sinh chỉ quen làm theo.
\textit{(She did not know what she liked, what she was good at, or what to choose, because her student years had trained her only to obey.)}

Một đứa trẻ quá ngoan đôi khi không cần thêm lời khen.
\textit{(An overly obedient child sometimes does not need more praise.)}

Nó cần được hỏi: ``Con thực sự nghĩ gì?''
\textit{(That child needs to be asked: ``What do you really think?'')}

\end{truyen}

\begin{baihoc}
Giáo dục không phải đào tạo ra người chỉ biết nghe lời. Giáo dục tốt phải giúp học sinh biết suy nghĩ, biết phản biện, và biết chịu trách nhiệm cho lựa chọn của mình.
\textit{(Education is not the training of compliant people. Good education should help students think, disagree respectfully, and own their choices.)}
\end{baihoc}

\begin{ghinhoanh}
\textbf{Short English to remember:}

Education is not the training of compliant people.

Good education should help students think, disagree respectfully, and own their choices.

\end{ghinhoanh}

\begin{tuhocnhanh}
\textbf{Cau hoi tu hoc / Self-study questions}

1. Van de chinh la gi? \textit{What is the main problem?}

2. Cach xu ly tot hon la gi? \textit{What is a better action?}

3. Mot cau tieng Anh em muon nho la cau nao? \textit{Which English sentence do you want to remember?}
\end{tuhocnhanh}
\ngancach

\section{Nói Chuyện Thật Với Cha Mẹ Là Một Kỹ Năng}
\begin{truyen}{Nói Chuyện Thật Với Cha Mẹ Là Một Kỹ Năng}{Speaking Honestly With Parents Is a Skill}
\chuhoa{K}{hông phải học sinh nào cũng có cha mẹ dễ nói chuyện.}
\textit{(Not every student has parents who are easy to talk to.)}

Có nhà vừa mở lời đã thành cãi nhau.
\textit{(In some homes, every conversation turns into conflict.)}

Nhưng im lặng mãi cũng không giải quyết được.
\textit{(But silence solves nothing either.)}

Có những điều chỉ thay đổi khi con cái học cách nói rõ, bình tĩnh, và có dẫn chứng.
\textit{(Some things change only when children learn to speak clearly, calmly, and with evidence.)}

Mai viết ra giấy ba điều: con đang mệt vì lịch học thêm nào, con muốn giảm gì, và con sẽ bù bằng cách nào.
\textit{(Mai wrote down three things: which extra classes were exhausting her, what she wanted reduced, and how she would make up for it.)}

Khi nói với mẹ, em không dùng câu trách móc như ``Mẹ ép con''. Em nói: ``Con đang quá tải. Nếu giữ lịch này, con học tệ hơn chứ không tốt hơn.''
\textit{(When she spoke to her mother, she did not use accusing lines like ``You force me.'' She said: ``I am overloaded. If this schedule stays, I will study worse, not better.'')}

Cuộc nói chuyện không hoàn hảo, nhưng đủ để mẹ em lùi một bước.
\textit{(The conversation was not perfect, but enough for her mother to step back.)}

Người lớn không phải lúc nào cũng thay đổi vì cảm xúc.
\textit{(Adults do not always change because of emotion.)}

Nhiều khi họ thay đổi vì cuối cùng cũng nghe được sự thật rõ ràng.
\textit{(Often they change because they finally hear the truth stated clearly.)}

\end{truyen}

\begin{baihoc}
Muốn được lắng nghe, học sinh cũng phải học cách nói có cấu trúc, không chỉ bùng nổ. Sự thật nói đúng cách có sức nặng hơn tiếng cãi vã.
\textit{(If you want to be heard, you must learn to speak with structure, not only with explosion. Truth said properly carries more weight than shouting.)}
\end{baihoc}

\begin{ghinhoanh}
\textbf{Short English to remember:}

If you want to be heard, you must learn to speak with structure, not only with explosion.

Truth said properly carries more weight than shouting.

\end{ghinhoanh}

\begin{tuhocnhanh}
\textbf{Cau hoi tu hoc / Self-study questions}

1. Van de chinh la gi? \textit{What is the main problem?}

2. Cach xu ly tot hon la gi? \textit{What is a better action?}

3. Mot cau tieng Anh em muon nho la cau nao? \textit{Which English sentence do you want to remember?}
\end{tuhocnhanh}
\ngancach

\section{Cha Muốn Con Nối Nghiệp}
\begin{truyen}{Cha Muốn Con Nối Nghiệp}{The Father Wants the Son to Continue the Line}
\chuhoa{N}{gày xưa trong một gia đình có nghề thuốc, người cha nhất quyết muốn con trai lớn nối nghiệp.}
\textit{(Long ago in a family of physicians, the father insisted that his eldest son continue the profession.)}

Nhưng cậu con lại thích đo đạc, dựng mô hình, và mê nghề xây dựng cầu đường.
\textit{(But the son loved measuring, making models, and was drawn to building bridges and roads.)}

Người cha giận dữ: ``Nhà ta bao đời cứu người, con bỏ đi là bất hiếu.''
\textit{(The father grew angry: ``Our family has healed people for generations. If you leave, you are unfilial.'')}

Bà nội chỉ nói một câu: "Cứu người bằng thuốc là cứu.
\textit{(The grandmother said only one thing: "Healing with medicine saves people.)}

Làm cây cầu để người ta qua sông an toàn cũng là cứu." Người cha im đi rất lâu.
\textit{(Building a bridge so people cross safely also saves people." The father fell silent for a long time.)}

\end{truyen}

\begin{baihoc}
Nhiều áp lực nghề nghiệp từ xưa đến nay đều đội lốt chữ hiếu. Nhưng hiếu không có nghĩa là sao chép nguyên xi cuộc đời cha mẹ.
\textit{(Across time, career pressure often wears the mask of filial duty. But duty does not mean copying a parent’s life exactly.)}
\end{baihoc}

\begin{ghinhoanh}
\textbf{Short English to remember:}

Across time, career pressure often wears the mask of filial duty.

But duty does not mean copying a parent’s life exactly.

\end{ghinhoanh}

\begin{tuhocnhanh}
\textbf{Cau hoi tu hoc / Self-study questions}

1. Van de chinh la gi? \textit{What is the main problem?}

2. Cach xu ly tot hon la gi? \textit{What is a better action?}

3. Mot cau tieng Anh em muon nho la cau nao? \textit{Which English sentence do you want to remember?}
\end{tuhocnhanh}
\ngancach

\section{Gánh Kỳ Vọng Của Con Cả}
\begin{truyen}{Gánh Kỳ Vọng Của Con Cả}{The Weight on the Eldest Child}
\chuhoa{T}{rong nhiều gia đình, đứa con đầu không chỉ là con.}
\textit{(In many families, the eldest child is not treated only as a child.)}

Nó còn bị coi như người phải mở đường cho cả nhà.
\textit{(That child is also treated as the one who must open the road for the whole family.)}

Hải học giỏi, nên mọi người đặt hết niềm tin vào em.
\textit{(Hai studied well, so the family placed all their hope on him.)}

Tin nhiều đến mức biến thành gánh.
\textit{(So much hope that it became a burden.)}

Mỗi lần em mệt, người lớn lại nhắc chuyện em là anh cả, phải làm gương, phải gánh.
\textit{(Whenever he tired, adults reminded him that he was the eldest, had to set an example, had to carry the load.)}

Không ai hỏi một đứa trẻ gánh lâu như vậy có đau vai không.
\textit{(No one asked whether a child carrying that weight for so long had aching shoulders.)}

\end{truyen}

\begin{baihoc}
Kỳ vọng có thể nâng người ta lên, nhưng quá tải thì nó nghiền người ta xuống. Con cả cũng là một đứa trẻ trước khi là biểu tượng.
\textit{(Expectation can lift someone up, but overload crushes them down. An eldest child is still a child before becoming a symbol.)}
\end{baihoc}

\begin{ghinhoanh}
\textbf{Short English to remember:}

Expectation can lift someone up, but overload crushes them down.

An eldest child is still a child before becoming a symbol.

\end{ghinhoanh}

\begin{tuhocnhanh}
\textbf{Cau hoi tu hoc / Self-study questions}

1. Van de chinh la gi? \textit{What is the main problem?}

2. Cach xu ly tot hon la gi? \textit{What is a better action?}

3. Mot cau tieng Anh em muon nho la cau nao? \textit{Which English sentence do you want to remember?}
\end{tuhocnhanh}
\ngancach

\section{Mẹ Bán Rau Và Bảng Điểm}
\begin{truyen}{Mẹ Bán Rau Và Bảng Điểm}{The Vegetable Seller Mother and the Report Card}
\chuhoa{M}{ột người mẹ bán rau ngoài chợ gửi hết hy vọng vào con gái.}
\textit{(A mother selling vegetables in the market placed all her hope in her daughter.)}

Bà tin chỉ có học thật giỏi mới thoát nghèo.
\textit{(She believed only excellent academic success could lift them from poverty.)}

Niềm tin đó không sai.
\textit{(That belief was not wrong.)}

Nhưng dần dần nó biến thành một câu vô hình treo trên đầu đứa bé: mày không được phép trượt.
\textit{(But gradually it became an invisible sentence hanging over the girl’s head: you are not allowed to fail.)}

Mỗi lần con được điểm thấp, người mẹ không đánh.
\textit{(Whenever her daughter scored low, the mother did not hit her.)}

Bà chỉ ngồi im và thở dài.
\textit{(She simply sat quietly and sighed.)}

Có những tiếng thở dài còn nặng hơn roi.
\textit{(Some sighs weigh heavier than a cane.)}

Đến khi cô giáo mời mẹ lên trường và nói: ``Chị đang làm con sợ học chứ không còn muốn học nữa,'' bà mới khóc.
\textit{(When the teacher invited the mother to school and said, ``You are making her fear study, not love it anymore,'' the mother finally cried.)}

\end{truyen}

\begin{baihoc}
Khát vọng cho con đổi đời là điều đáng trọng. Nhưng nếu không tỉnh táo, hy vọng của người lớn sẽ biến thành nỗi sợ của trẻ nhỏ.
\textit{(A parent’s hope for a better life is honorable. But without awareness, adult hope turns into a child’s fear.)}
\end{baihoc}

\begin{ghinhoanh}
\textbf{Short English to remember:}

A parent’s hope for a better life is honorable.

But without awareness, adult hope turns into a child’s fear.

\end{ghinhoanh}

\begin{tuhocnhanh}
\textbf{Cau hoi tu hoc / Self-study questions}

1. Van de chinh la gi? \textit{What is the main problem?}

2. Cach xu ly tot hon la gi? \textit{What is a better action?}

3. Mot cau tieng Anh em muon nho la cau nao? \textit{Which English sentence do you want to remember?}
\end{tuhocnhanh}
\ngancach

\section{Nhật Ký Không Dám Đưa Cho Mẹ}
\begin{truyen}{Nhật Ký Không Dám Đưa Cho Mẹ}{The Diary She Dared Not Show Her Mother}
\chuhoa{M}{ột học sinh viết trong nhật ký: "Con không sợ học khó.}
\textit{(A student wrote in her diary: "I am not afraid of difficult study.)}

Con sợ làm mẹ buồn."
\textit{(I am afraid of making my mother sad.")}

Dòng đó ngắn, nhưng nói hết tình cảnh của rất nhiều đứa trẻ học giỏi mà mệt mỏi.
\textit{(The line was short, but it captured the condition of many high-performing yet exhausted children.)}

Các em không chỉ đối mặt kiến thức.
\textit{(They are not facing knowledge alone.)}

Các em còn phải gánh trạng thái cảm xúc của người lớn.
\textit{(They are also carrying adult emotional states.)}

Khi con cái học để giữ hòa khí gia đình hơn là học cho tương lai của chính mình, việc học rất dễ hóa thành gông.
\textit{(When children study more to preserve family peace than for their own future, study easily becomes a yoke.)}

\end{truyen}

\begin{baihoc}
Học sinh không nên phải làm nhiệm vụ giữ cân bằng cảm xúc cho cha mẹ. Đó là gánh quá lớn với một người còn đang lớn.
\textit{(Students should not have to stabilize their parents’ emotions. That burden is too large for someone still growing up.)}
\end{baihoc}

\begin{ghinhoanh}
\textbf{Short English to remember:}

Students should not have to stabilize their parents’ emotions.

That burden is too large for someone still growing up.

\end{ghinhoanh}

\begin{tuhocnhanh}
\textbf{Cau hoi tu hoc / Self-study questions}

1. Van de chinh la gi? \textit{What is the main problem?}

2. Cach xu ly tot hon la gi? \textit{What is a better action?}

3. Mot cau tieng Anh em muon nho la cau nao? \textit{Which English sentence do you want to remember?}
\end{tuhocnhanh}
\ngancach

\section{Cuộc Nói Chuyện Sau Bữa Cơm}
\begin{truyen}{Cuộc Nói Chuyện Sau Bữa Cơm}{The Conversation After Dinner}
\chuhoa{S}{au nhiều tháng im lặng, một cậu con trai đặt bảng điểm xuống bàn và nói: ``Con cần bố nghe hết trước khi đánh giá.''}
\textit{(After months of silence, a son placed his report card on the table and said, ``I need you to hear me through before judging.'')}

Cậu kể mình đang học vì sợ, ngủ không ngon, và không còn thấy vui với bất kỳ môn nào nữa.
\textit{(He explained that he was studying out of fear, sleeping badly, and no longer feeling joy in any subject.)}

Người bố ban đầu khó chịu.
\textit{(The father was irritated at first.)}

Nhưng khi nghe đến câu: ``Con không dám kể gì với bố vì lúc nào bố cũng chỉ hỏi điểm,'' ông khựng lại.
\textit{(But when he heard, ``I do not dare tell you anything because you only ever ask about scores,'' he stopped.)}

Không phải gia đình nào cũng đổi ngay.
\textit{(Not every family changes immediately.)}

Nhưng nhiều thay đổi bắt đầu từ một bữa cơm có người chịu nói thật.
\textit{(But many changes begin with one dinner where someone dares to tell the truth.)}

\end{truyen}

\begin{baihoc}
Nhiều áp lực tồn tại dai vì trong nhà không ai nói thật. Giao tiếp vụng vẫn tốt hơn im lặng kéo dài.
\textit{(Many pressures survive because no one in the family speaks honestly. Awkward communication is still better than prolonged silence.)}
\end{baihoc}

\begin{ghinhoanh}
\textbf{Short English to remember:}

Many pressures survive because no one in the family speaks honestly.

Awkward communication is still better than prolonged silence.

\end{ghinhoanh}

\begin{tuhocnhanh}
\textbf{Cau hoi tu hoc / Self-study questions}

1. Van de chinh la gi? \textit{What is the main problem?}

2. Cach xu ly tot hon la gi? \textit{What is a better action?}

3. Mot cau tieng Anh em muon nho la cau nao? \textit{Which English sentence do you want to remember?}
\end{tuhocnhanh}
\ngancach

\section{Người Lớn Cũng Cần Học Cách Xin Lỗi}
\begin{truyen}{Người Lớn Cũng Cần Học Cách Xin Lỗi}{Adults Also Need to Learn to Apologize}
\chuhoa{M}{ột ông bố từng ép con trai học thêm kín tuần.}
\textit{(A father once forced his son into a week packed with extra classes.)}

Đến khi thấy con hoảng loạn trước mỗi bài kiểm tra, ông mới nhận ra mình đã đi quá.
\textit{(Only when he saw the boy panic before every test did he realize he had gone too far.)}

Tối đó ông nói một câu hiếm: "Bố xin lỗi.
\textit{(That night he said something rare: "I am sorry.)}

Bố tưởng thúc con là giúp con."
\textit{(I thought pushing you was helping you.")}

Đứa con ngồi im rồi khóc.
\textit{(The child sat silently and cried.)}

Không phải vì câu xin lỗi xóa hết mọi thứ, mà vì lần đầu tiên em thấy người lớn chịu nhận phần sai của mình.
\textit{(Not because the apology erased everything, but because for the first time he saw an adult admit fault.)}

Nhiều gia đình không thiếu tình thương.
\textit{(Many families do not lack love.)}

Họ chỉ thiếu khả năng sửa sai sau khi đã thương sai cách.
\textit{(They lack the ability to correct themselves after loving in the wrong way.)}

\end{truyen}

\begin{baihoc}
Muốn giáo dục con cái tốt, người lớn cũng phải học. Một lời xin lỗi đúng lúc có thể mở ra không khí mới cho cả nhà.
\textit{(To educate children well, adults must also learn. A timely apology can open a new atmosphere for the whole family.)}
\end{baihoc}

\begin{ghinhoanh}
\textbf{Short English to remember:}

To educate children well, adults must also learn.

A timely apology can open a new atmosphere for the whole family.

\end{ghinhoanh}

\begin{tuhocnhanh}
\textbf{Cau hoi tu hoc / Self-study questions}

1. Van de chinh la gi? \textit{What is the main problem?}

2. Cach xu ly tot hon la gi? \textit{What is a better action?}

3. Mot cau tieng Anh em muon nho la cau nao? \textit{Which English sentence do you want to remember?}
\end{tuhocnhanh}
